\section{Motivation and Design Principles}

Standard self-attention computes a dense $n \times n$ interaction matrix, yet empirical
studies \cite{clark2019does, kovaleva2019revealing} have shown that attention maps in
trained transformers are highly sparse: only a small fraction of token pairs receive
non-negligible attention weights. This observation suggests that explicit sparsity
constraints, if applied intelligently, may reduce computation without sacrificing
expressiveness.

The key challenge is that the sparsity pattern is \textit{context-dependent}:
which tokens a given query should attend to depends on the input itself. Fixed patterns
— such as local windows — cannot capture long-range dependencies that span arbitrary
distances. Our Sparse Contextual Attention mechanism addresses this by learning to
predict, for each query, a small set of relevant keys directly from the input
representations.
